\documentclass[../main.tex]{subfiles}

\begin{document}
\section*{AI Acknowledgement}
\addcontentsline{toc}{section}{AI Acknowledgement}

\subsection*{What generative AI tools have you used to complete A2?}
\noindent
We used \textbf{Anthropic Claude (Opus 4.7, 1M-context)} through the
\emph{Claude Code} command-line agent as a coding-and-writing assistant.
The model and tool ran on the local development machine and was the only
generative AI tool used.

\subsection*{How have you used generative AI tools in A2?}
\noindent
The assistant was used in three roles:
\begin{itemize}
    \item \textbf{Notebook engineering.} Refactoring Section~1 of the
    notebook into clearly delimited subsections (data loading, class
    definitions, sample sizes, sample images, pixel statistics,
    pre-processing, and a stratified train/validation split), drafting the
    \texttt{make\_svm}/\texttt{make\_mlp}/\texttt{make\_cnn} factory
    functions, and writing the grid-search and final-evaluation cells in
    Sections~3 and~4.
    \item \textbf{Report drafting.} Producing first drafts of the
    Introduction, Data, Methods, Results~\&~Discussion and Conclusion
    sections, and assembling the LaTeX tables that summarise the
    hyper-parameter grids and final test performance.
    \item \textbf{Tooling and verification.} Setting up the
    \texttt{tensorflow-metal} environment so that training ran on the
    Apple~M3 GPU, and running quick sanity checks (shape audits,
    anonymisation scans).
\end{itemize}

\noindent
All AI-suggested code and prose were read, validated and edited by the
group members. Every numerical result reported in this document was
obtained by executing the (group-reviewed) code on our own hardware; no
results, citations or claims were taken from the assistant without a
human in the loop.

\end{document}
